\chapter[Strategi Bisnis dan Rumusan Masalah]{Strategi Bisnis dan Rumusan
Masalah: Menyelaraskan Arah Organisasi dengan Pertanyaan Penelitian}
\label{ch:strategi-bisnis}



\section*{Tujuan Pembelajaran}
\begin{itemize}
  \item Menjelaskan hubungan antara strategi organisasi, model bisnis,
        kapabilitas utama, dan kebutuhan arsitektur enterprise dengan bahasa
        manajerial.
  \item Menggunakan kanvas model bisnis dan konsep \emph{operating model}
        untuk membaca arah strategis organisasi studi kasus.
  \item Mengidentifikasi ketidakselarasan antara strategi bisnis dan keputusan
        teknologi sebagai sumber masalah penelitian.
  \item Merumuskan rumusan masalah, pertanyaan penelitian (RQ), dan kontribusi
        awal yang layak dikembangkan menjadi makalah publikasi individual.
\end{itemize}


\section{Ringkasan Awal Bab}

Bab ini membahas cara mengubah topik awal menjadi rumusan masalah, pertanyaan
penelitian, dan kontribusi makalah. Pembahasan dimulai dari dilema organisasi:
strategi digital sering terdengar jelas di tingkat slogan, tetapi tidak selalu
selaras dengan proses, data, aplikasi, teknologi, API, dan tata kelola yang
dimiliki organisasi. Untuk membaca ketegangan tersebut, bab ini menggunakan
kanvas model bisnis dan konsep \emph{operating model} sebagai alat analisis
yang ramah bagi pembaca non-IT. Melalui studi kasus Bank Wanua, pembaca
melihat bagaimana strategi menjadi mitra digital UMKM dapat diterjemahkan
menjadi daftar ketidakselarasan, lalu diubah menjadi rumusan masalah, RQ, dan
pernyataan kontribusi. Keluaran bab ini memperkuat bagian Pendahuluan makalah,
terutama latar belakang masalah, RQ, dan kontribusi awal.

\section{Pendahuluan}

Bab sebelumnya membantu pembaca memilih organisasi studi kasus, menulis
pernyataan topik awal, dan mengenali target publikasi yang realistis. Bab ini
melangkah satu tahap lebih jauh: topik awal tersebut harus diubah menjadi
masalah penelitian yang tajam. Dalam makalah publikasi, masalah yang tajam
lebih penting daripada topik yang terdengar menarik. Topik menyebut area
umum, sedangkan masalah menjelaskan ketegangan yang layak dianalisis, bukti
yang perlu dikumpulkan, dan kontribusi yang dapat diberikan.

Bagi pembaca non-IT, strategi bisnis menjadi titik mulai yang sehat. Banyak
inisiatif digital gagal bukan karena teknologi yang dipilih terlalu lemah,
melainkan karena arah bisnis, model operasi, dan prioritas teknologi tidak
dibaca sebagai satu kesatuan. Organisasi ingin bergerak cepat, tetapi datanya
tersebar. Organisasi ingin membuka layanan digital, tetapi proses internal
masih manual. Organisasi ingin membangun platform, tetapi belum jelas siapa
pengguna utama dan nilai apa yang dipertukarkan. Dalam bahasa Enterprise
Architecture, persoalan seperti ini adalah persoalan keselarasan antara
strategi dan kapabilitas.

Rujukan klasik Enterprise Architecture menempatkan arsitektur sebagai fondasi
pelaksanaan strategi bisnis, bukan sekadar urusan diagram teknologi
\cite{ross2006enterprise}. Pandangan yang lebih baru juga menekankan bahwa
organisasi digital membutuhkan kapabilitas yang dirancang secara sadar agar
strategi dapat dieksekusi berulang, aman, dan terukur \cite{ross2019designed}.
Karena itu, Bab 2 mengajak pembaca membaca strategi organisasi dengan dua alat
sederhana: kanvas model bisnis dan konsep \emph{operating model}. Keduanya
dipakai untuk menyusun rumusan masalah, RQ, dan pernyataan kontribusi awal.

\section{Skenario Pembuka: Ketika Strategi Digital Tidak Menyatu}

Bank Wanua, yang diperkenalkan pada Bab~\ref{ch:orientasi}, kini sudah
menetapkan ambisi strategis baru. Dalam tiga tahun, bank daerah ini ingin
menjadi mitra keuangan digital utama bagi UMKM lokal. Direksi membayangkan
aplikasi mobile yang lebih baik, layanan pembayaran yang terhubung dengan
ekosistem QRIS, API untuk mitra koperasi dan marketplace lokal, serta program
analitik sederhana untuk membaca kebutuhan nasabah UMKM.

Namun setelah beberapa bulan berjalan, tanda-tanda ketidakselarasan mulai
muncul. Tim pemasaran menjanjikan pembukaan rekening digital yang cepat, tetapi
tim operasional masih membutuhkan verifikasi manual berlapis. Tim TI diminta
membangun API untuk mitra, tetapi daftar layanan prioritas belum disepakati.
Divisi kepatuhan meminta kontrol data yang ketat, sementara unit bisnis ingin
kampanye yang lebih personal. Kantor cabang ingin mempertahankan relasi tatap
muka dengan nasabah, tetapi kantor pusat menekan efisiensi melalui kanal
digital.

Masalah Bank Wanua bukan sekadar ``belum punya aplikasi yang bagus''. Masalah
yang lebih mendasar adalah belum jelasnya hubungan antara strategi, model
bisnis, model operasi, dan keputusan arsitektur. Jika mahasiswa langsung
menulis makalah tentang ``perancangan aplikasi mobile Bank Wanua'', naskahnya
mungkin terlalu sempit. Jika mahasiswa menulis tentang ``transformasi digital
Bank Wanua'', naskahnya mungkin terlalu luas. Tantangannya adalah menemukan
rumusan masalah yang berada di tengah: cukup sempit untuk dianalisis dalam
satu semester, tetapi cukup strategis untuk memberi kontribusi akademik dan
praktis.

\subsection*{Pertanyaan Pemandu Bab Ini}

\begin{enumerate}
  \item Strategi bisnis apa yang sebenarnya sedang dikejar organisasi studi
        kasus?
  \item Bagaimana model bisnis organisasi menghasilkan nilai bagi pelanggan,
        mitra, dan pemangku kepentingan lain?
  \item Model operasi seperti apa yang dibutuhkan agar strategi tersebut dapat
        dijalankan secara konsisten?
  \item Di mana letak ketidakselarasan antara strategi, proses, data, aplikasi,
        teknologi, API, dan tata kelola?
  \item Bagaimana ketidakselarasan itu diubah menjadi rumusan masalah, RQ, dan
        kontribusi awal makalah?
\end{enumerate}

\section{Mengapa Topik Ini Penting bagi Pembaca Non-IT}

Pembaca non-IT tidak harus menjadi arsitek teknis untuk dapat membaca masalah
strategi digital. Yang dibutuhkan adalah kemampuan melihat pola: apa tujuan
organisasi, siapa pengguna yang dilayani, nilai apa yang dijanjikan, proses apa
yang harus berubah, data apa yang dibutuhkan, dan keputusan apa yang berisiko
jika tidak dikelola. Kemampuan ini sering justru berada dekat dengan peran
manajerial, karena manajer melihat dampak keputusan teknologi pada pelanggan,
pegawai, anggaran, regulasi, dan reputasi.

Dalam organisasi nyata, pertanyaan teknis hampir selalu membawa konsekuensi
manajerial. Membuka API bukan hanya pekerjaan membuat endpoint; ia mengubah
cara organisasi berhubungan dengan mitra. Memusatkan data pelanggan bukan
hanya pekerjaan integrasi; ia memunculkan pertanyaan tentang kepemilikan data,
izin pemrosesan, dan perlindungan privasi. Mengganti aplikasi lama bukan hanya
pekerjaan pengadaan; ia mengganggu kebiasaan kerja, proses persetujuan, dan
indikator kinerja.

Karena itu, rumusan masalah dalam makalah Enterprise Architecture tidak cukup
berbunyi ``sistem belum terintegrasi''. Rumusan yang lebih kuat menjelaskan
mengapa integrasi itu penting bagi strategi, siapa yang terkena dampaknya, apa
bukti ketidakselarasannya, dan bagaimana analisis arsitektur dapat membantu
pengambil keputusan. Bab ini memberi kerangka untuk melakukan penyempitan
tersebut.

\section{Konsep Inti: Strategi, Model Bisnis, dan RQ}

Empat konsep inti digunakan dalam bab ini. Pertama, \textbf{strategi bisnis}
menjelaskan pilihan organisasi tentang posisi, pelanggan, nilai, dan prioritas.
Kedua, \textbf{model bisnis} menjelaskan bagaimana organisasi menciptakan,
menyampaikan, dan menangkap nilai. Ketiga, \textbf{operating model} menjelaskan
tingkat standardisasi proses dan integrasi data yang dibutuhkan agar strategi
dapat dijalankan \cite{ross2006enterprise}. Keempat, \textbf{rumusan masalah
dan RQ} menerjemahkan ketegangan strategis menjadi pertanyaan akademik yang
dapat dijawab melalui analisis.

\subsection{Strategi Bisnis: Pilihan, Bukan Slogan}

Strategi sering ditulis sebagai slogan: ``menjadi pemimpin digital'', ``memberi
layanan terbaik'', atau ``mendorong inovasi''. Slogan semacam itu tidak salah,
tetapi belum cukup untuk menjadi dasar arsitektur. Strategi yang dapat
diterjemahkan ke dalam arsitektur harus memperlihatkan pilihan. Pilihan itu
meliputi pelanggan mana yang diprioritaskan, layanan mana yang menjadi
pembeda, proses mana yang perlu distandardisasi, dan risiko mana yang bersedia
ditanggung.

Dalam konteks Bank Wanua, strategi ``menjadi mitra digital UMKM'' baru
bermakna jika diterjemahkan ke pertanyaan yang lebih konkret. Apakah bank
ingin memperluas basis nasabah baru, memperdalam hubungan dengan nasabah
lama, menjadi penyedia API untuk mitra, atau menjadi bagian dari ekosistem
platform yang lebih besar? Setiap pilihan membawa konsekuensi arsitektur yang
berbeda. Strategi akuisisi nasabah membutuhkan kanal digital dan verifikasi
identitas. Strategi kemitraan API membutuhkan kontrak layanan, keamanan, dan
tata kelola mitra. Strategi efisiensi cabang membutuhkan perubahan proses
operasi.

\subsection{Kanvas Model Bisnis: Membaca Cara Organisasi Menghasilkan Nilai}

Kanvas model bisnis membantu pembaca memetakan organisasi dalam sembilan
komponen yang mudah dibaca. Untuk kebutuhan bab ini, kanvas tidak dipakai
sebagai latihan kewirausahaan, melainkan sebagai alat untuk menghubungkan
strategi dengan kapabilitas arsitektur.

\begin{table}[htbp]
\centering
\scriptsize
\caption{Pertanyaan panduan kanvas model bisnis untuk analisis arsitektur.}
\label{tab:bmc-pertanyaan}
\begin{tabularx}{\textwidth}{@{}p{0.26\textwidth} X@{}}
\hline
\textbf{Komponen} & \textbf{Pertanyaan Arsitektural} \\
\hline
Segmen pelanggan & Kelompok pengguna mana yang paling penting, dan data apa
yang dibutuhkan untuk memahami kebutuhan mereka? \\
Proposisi nilai & Nilai apa yang dijanjikan, dan kapabilitas apa yang harus
ada agar janji itu dapat ditepati? \\
Kanal & Kanal fisik dan digital apa yang dipakai, dan bagaimana kanal tersebut
perlu terintegrasi? \\
Relasi pelanggan & Hubungan seperti apa yang dibangun: personal, otomatis,
komunitas, atau berbasis mitra? \\
Sumber pendapatan & Nilai ekonomi apa yang dihasilkan, dan metrik apa yang
dapat menunjukkan keberhasilannya? \\
Aktivitas kunci & Proses apa yang menentukan keberhasilan strategi? \\
Sumber daya kunci & Data, aplikasi, SDM, infrastruktur, dan mitra apa yang
menjadi sumber kemampuan utama? \\
Mitra kunci & Pihak luar mana yang perlu terhubung, dan aturan integrasi apa
yang dibutuhkan? \\
Struktur biaya & Biaya apa yang paling besar, dan keputusan arsitektur apa
yang dapat mengendalikannya? \\
\hline
\end{tabularx}
\end{table}

Kekuatan kanvas ini adalah kesederhanaannya. Pembaca dapat menggunakannya
untuk melihat apakah masalah digital yang dipilih benar-benar terhubung dengan
nilai organisasi. Jika sebuah inisiatif teknologi tidak jelas hubungannya
dengan pelanggan, proposisi nilai, aktivitas kunci, atau mitra, maka inisiatif
itu perlu dipertanyakan sebelum masuk ke rancangan arsitektur.

\subsection{Operating Model: Tingkat Integrasi dan Standardisasi}

\emph{Operating model} menjelaskan bagaimana organisasi menjalankan strategi
melalui proses dan data. Dalam rujukan Enterprise Architecture, dua pertanyaan
kunci digunakan untuk membaca model operasi: seberapa besar proses perlu
distandardisasi, dan seberapa besar data perlu diintegrasikan
\cite{ross2006enterprise}. Dari dua pertanyaan ini muncul empat pola umum.

\begin{figure}[htbp]
\centering
\scalebox{0.95}{
\begin{tikzpicture}[
  font=\small,
  axis/.style={-Latex, thick},
  cell/.style={draw, rounded corners=1.5mm, align=center,
    minimum width=46mm, minimum height=24mm},
  note/.style={font=\footnotesize\mdseries, align=center, text width=42mm}
]
  \draw[axis] (0,0) -- (10,0) node[right] {\textbf{Integrasi data tinggi}};
  \draw[axis] (0,0) -- (0,6) node[above, align=center] {\textbf{Standardisasi}\\\textbf{proses tinggi}};

  \node[cell, fill=gray!12] (div) at (2.5,1.5)
    {\textbf{Diversification}\\\footnotesize unit otonom};
  \node[cell, fill=blue!12] (coord) at (7.5,1.5)
    {\textbf{Coordination}\\\footnotesize data terintegrasi};
  \node[cell, fill=orange!18] (rep) at (2.5,4.5)
    {\textbf{Replication}\\\footnotesize proses seragam};
  \node[cell, fill=praditagreen!20] (uni) at (7.5,4.5)
    {\textbf{Unification}\\\footnotesize proses dan data seragam};

  \node[note, below=2mm of div] {proses berbeda, data tidak banyak dibagi};
  \node[note, below=2mm of coord] {proses boleh berbeda, data perlu menyatu};
  \node[note, above=2mm of rep] {proses sama, data lokal};
  \node[note, above=2mm of uni] {proses sama, data bersama};
\end{tikzpicture}
}
\caption{Empat pola operating model berdasarkan kebutuhan standardisasi proses
dan integrasi data.}
\label{fig:operating-model}
\end{figure}

Pola \emph{Diversification} cocok ketika unit bisnis relatif mandiri. Pola
\emph{Coordination} cocok ketika unit dapat bekerja berbeda, tetapi perlu
berbagi data pelanggan atau transaksi. Pola \emph{Replication} cocok ketika
organisasi ingin menyalin proses yang sama di banyak lokasi. Pola
\emph{Unification} cocok ketika proses dan data perlu sangat seragam, misalnya
untuk kepatuhan, layanan nasional, atau efisiensi besar-besaran.

Pemilihan operating model bukan sekadar label. Ia membantu mahasiswa
menjelaskan mengapa sebuah masalah arsitektur muncul. Jika strategi organisasi
membutuhkan koordinasi data pelanggan lintas unit, tetapi setiap unit
menyimpan data dengan definisi berbeda, maka ketidakselarasan itu dapat menjadi
sumber rumusan masalah. Jika organisasi menginginkan layanan digital yang sama
di seluruh cabang, tetapi proses cabang sangat berbeda, maka masalahnya bukan
hanya aplikasi, melainkan standardisasi proses.

\subsection{Rumusan Masalah, RQ, dan Kontribusi}

Rumusan masalah adalah pernyataan tentang ketegangan yang layak dianalisis.
Pertanyaan penelitian atau RQ adalah bentuk tanya yang membatasi arah
analisis. Kontribusi adalah janji tentang nilai yang akan diberikan makalah
kepada pembaca akademik atau praktisi.

Hubungan ketiganya dapat diringkas sebagai berikut:

\begin{itemize}
  \item \textbf{Rumusan masalah} menjawab: apa ketegangan utama yang terjadi?
  \item \textbf{RQ} menjawab: pertanyaan apa yang akan dijawab makalah?
  \item \textbf{Kontribusi} menjawab: apa manfaat jawaban tersebut bagi
        pengetahuan atau praktik?
\end{itemize}

Contoh rumusan yang masih lemah adalah: ``Bank Wanua perlu melakukan
transformasi digital.'' Pernyataan ini terlalu luas dan belum menunjukkan
ketegangan. Rumusan yang lebih baik: ``Strategi Bank Wanua untuk menjadi mitra
digital UMKM belum didukung oleh model operasi yang jelas, terutama pada
integrasi data nasabah, prioritas layanan API, dan pembagian tanggung jawab
antara kantor pusat, cabang, dan mitra.'' Rumusan kedua lebih siap diubah
menjadi RQ.

\section{Glosarium Mini}

\begin{itemize}
  \item \textbf{Strategi bisnis:} pilihan organisasi tentang pelanggan,
        proposisi nilai, prioritas, dan cara bersaing atau melayani.
  \item \textbf{Model bisnis:} cara organisasi menciptakan, menyampaikan, dan
        menangkap nilai melalui pelanggan, kanal, aktivitas, sumber daya, dan
        mitra.
  \item \textbf{Operating model:} pola pelaksanaan strategi yang menjelaskan
        kebutuhan standardisasi proses dan integrasi data.
  \item \textbf{Ketidakselarasan bisnis--TI:} jarak antara tujuan bisnis dan
        kapabilitas teknologi, proses, data, aplikasi, atau tata kelola yang
        tersedia.
  \item \textbf{Rumusan masalah:} penjelasan ringkas tentang ketegangan utama
        yang layak dianalisis dalam makalah.
  \item \textbf{Research question (RQ):} pertanyaan penelitian yang membatasi
        arah analisis dan harus dijawab secara eksplisit dalam makalah.
  \item \textbf{Pernyataan kontribusi:} penjelasan tentang nilai akademik atau
        praktis yang diberikan makalah.
\end{itemize}

\section{Kerangka Berpikir: Dari Strategi ke Pertanyaan Penelitian}

Bab ini menggunakan alur berpikir yang sederhana. Pembaca mulai dari strategi
organisasi, memetakan model bisnis, membaca operating model, menemukan
ketidakselarasan, lalu merumuskan masalah, RQ, dan kontribusi awal. Alur ini
ditunjukkan pada Gambar~\ref{fig:strategi-ke-rq}.

\begin{figure}[htbp]
\centering
\scalebox{0.92}{
\begin{tikzpicture}[
  font=\small\bfseries,
  node distance=9mm,
  flowstep/.style={draw, rounded corners=1.5mm, align=center,
    minimum width=28mm, minimum height=14mm, text width=27mm},
  arrow/.style={-Latex, very thick}
]
  \node[flowstep, fill=praditagreen!20] (s) {Strategi\\organisasi};
  \node[flowstep, fill=orange!18, right=of s] (bmc) {Model\\bisnis};
  \node[flowstep, fill=blue!12, right=of bmc] (op) {Operating\\model};
  \node[flowstep, fill=gray!16, below=of bmc] (mis) {Ketidak-\\selarasan};
  \node[flowstep, fill=orange!20, left=of mis] (prob) {Rumusan\\masalah};
  \node[flowstep, fill=praditagreen!22, right=of mis] (rq) {RQ dan\\kontribusi};

  \draw[arrow] (s) -- (bmc);
  \draw[arrow] (bmc) -- (op);
  \draw[arrow] (op.south) -- (mis.north east);
  \draw[arrow] (bmc.south) -- (mis.north);
  \draw[arrow] (mis) -- (prob);
  \draw[arrow] (mis) -- (rq);
  \draw[arrow] (prob) -- (rq);

  \node[align=center, font=\footnotesize\mdseries, text width=96mm,
        below=12mm of mis]
    {Artefak utama bab ini: ringkasan strategi, kanvas model bisnis,
     klasifikasi operating model, daftar ketidakselarasan, draf rumusan
     masalah, RQ, dan kontribusi awal.};
\end{tikzpicture}
}
\caption{Alur berpikir Bab 2 dari strategi organisasi menuju pertanyaan
penelitian dan kontribusi makalah.}
\label{fig:strategi-ke-rq}
\end{figure}

Alur ini membantu mahasiswa menjaga fokus. Tanpa alur seperti ini, mahasiswa
sering melompat dari cerita organisasi ke rekomendasi teknologi. Padahal
makalah yang baik perlu memperlihatkan jejak penalaran: mengapa masalah itu
dipilih, mengapa penting, mengapa kerangka analisis tertentu dipakai, dan
mengapa rekomendasi yang muncul dapat dipertanggungjawabkan.

\section{Konteks Indonesia: Strategi Digital dalam Realitas Lokal}

Strategi digital organisasi Indonesia dibentuk oleh beberapa kondisi yang
perlu diperhatikan sejak awal. Pertama, kematangan digital antarorganisasi dan
antarwilayah tidak merata. Perusahaan besar di kota utama mungkin sudah
terbiasa dengan layanan berbasis API, sedangkan koperasi, BPR, rumah sakit
daerah, atau pemerintah kabupaten masih bergulat dengan integrasi dasar dan
kualitas data. Ketidakmerataan ini memengaruhi kelayakan rekomendasi.

Kedua, regulasi memberi batas dan arah. UU Pelindungan Data Pribadi menuntut
organisasi lebih berhati-hati dalam mengumpulkan, memproses, dan membagikan
data pribadi \cite{uupdp2022}. PP No.~71/2019 mengatur penyelenggaraan sistem
dan transaksi elektronik, termasuk tanggung jawab penyelenggara sistem
\cite{pp712019}. Di sektor pembayaran, SNAP BI dari Bank Indonesia memberi
standar integrasi API yang penting bagi organisasi keuangan dan fintech
\cite{snapbi}. Strategi yang tidak memperhitungkan regulasi dapat terlihat
menarik di atas kertas tetapi rapuh saat dijalankan.

Ketiga, banyak organisasi Indonesia bekerja dalam struktur yang hierarkis.
Keputusan digital sering tersentralisasi di kantor pusat, sementara
implementasi terjadi di cabang, unit layanan, atau mitra lapangan. Situasi ini
membuat operating model menjadi penting. Strategi digital yang menuntut
standardisasi tinggi membutuhkan mekanisme pelaksanaan yang berbeda dari
strategi yang memberi otonomi besar kepada unit lokal.

Keempat, biaya dan kapasitas SDM sering menjadi batas nyata. Mahasiswa tidak
boleh menulis rekomendasi seperti ``bangun platform nasional berbasis cloud
dan microservices'' tanpa menjelaskan prioritas, tahapan, risiko, dan
kesiapan organisasi. Dalam makalah, kepekaan terhadap konteks Indonesia bukan
hiasan; ia menentukan apakah analisis dapat dipercaya.

\section{Studi Kasus Utama: Menajamkan Masalah Bank Wanua}

Untuk melihat cara kerja bab ini, kita kembali ke Bank Wanua. Strategi awal
bank adalah menjadi mitra digital UMKM daerah. Strategi ini dapat dipetakan
ke dalam beberapa elemen kanvas model bisnis sebagaimana Tabel~\ref{tab:bank-wanua-bmc}.

\begin{table}[htbp]
\centering
\scriptsize
\caption{Ringkasan model bisnis Bank Wanua untuk strategi UMKM digital.}
\label{tab:bank-wanua-bmc}
\begin{tabularx}{\textwidth}{@{}p{0.24\textwidth} X X@{}}
\hline
\textbf{Elemen} & \textbf{Kondisi Saat Ini} & \textbf{Arah Strategis} \\
\hline
Segmen pelanggan & Nasabah ritel lokal, pegawai daerah, pelaku UMKM lama. &
UMKM yang membutuhkan layanan pembayaran, pembiayaan, dan pencatatan transaksi
digital. \\
Proposisi nilai & Kedekatan lokal dan layanan cabang. & Layanan keuangan
terintegrasi untuk UMKM: rekening, pembayaran, pembiayaan, dan akses mitra. \\
Kanal & Cabang, agen, ATM, aplikasi mobile dasar. & Mobile banking, API mitra,
QRIS, dan kanal cabang sebagai pendamping digital. \\
Aktivitas kunci & Pembukaan rekening, transaksi, kredit, layanan cabang. &
Onboarding digital UMKM, integrasi pembayaran, analitik risiko sederhana,
tata kelola mitra API. \\
Mitra kunci & Pemerintah daerah, koperasi, komunitas UMKM. & Fintech,
marketplace lokal, koperasi, penyedia layanan pembayaran, pemerintah daerah. \\
\hline
\end{tabularx}
\end{table}

Dari tabel tersebut terlihat bahwa strategi baru menuntut koordinasi data dan
standardisasi proses yang lebih tinggi. Bank Wanua tidak cukup hanya menambah
fitur aplikasi. Ia perlu memastikan data nasabah UMKM dapat dibaca lintas
kanal, proses pembukaan rekening dan verifikasi risiko cukup konsisten, dan
aturan berbagi data dengan mitra dipahami sejak awal. Dengan kata lain, Bank
Wanua bergerak dari model operasi yang cenderung cabang-sentris menuju model
yang lebih terkoordinasi.

Ketidakselarasan yang muncul dapat diringkas sebagai berikut:

\begin{enumerate}
  \item \textbf{Strategi pelanggan tidak selaras dengan data.} Bank ingin
        memahami UMKM secara lebih baik, tetapi data nasabah, transaksi, dan
        interaksi cabang belum menjadi satu pandangan yang konsisten.
  \item \textbf{Strategi kemitraan tidak selaras dengan tata kelola API.} Bank
        ingin terhubung dengan mitra, tetapi belum memiliki prioritas layanan
        API, kriteria mitra, dan mekanisme pengawasan.
  \item \textbf{Strategi efisiensi tidak selaras dengan proses cabang.} Bank
        ingin mempercepat layanan, tetapi proses verifikasi dan persetujuan
        masih berbeda antarcabang.
  \item \textbf{Strategi personalisasi tidak selaras dengan kepatuhan data.}
        Bank ingin memberi penawaran yang lebih personal, tetapi belum
        menjelaskan batas penggunaan data pribadi dan persetujuan nasabah.
\end{enumerate}

Dari sini, rumusan masalah awal dapat ditulis:

\begin{quote}
Strategi Bank Wanua untuk menjadi mitra digital UMKM daerah belum didukung
oleh model operasi dan arah arsitektur enterprise yang jelas. Ketidakselarasan
terlihat pada integrasi data nasabah UMKM, standardisasi proses lintas cabang,
prioritas layanan API untuk mitra, dan tata kelola penggunaan data pribadi.
Kondisi ini menyulitkan bank untuk menerjemahkan ambisi digital menjadi
kapabilitas yang dapat dijalankan secara bertahap dan sesuai konteks regulasi
Indonesia.
\end{quote}

Rumusan tersebut dapat diubah menjadi RQ awal:

\begin{quote}
Bagaimana pendekatan Enterprise Architecture berbasis TOGAF-lite dan BDAT
dapat digunakan untuk merumuskan arah arsitektur target yang menyelaraskan
strategi Bank Wanua sebagai mitra digital UMKM dengan kebutuhan integrasi data,
standardisasi proses, layanan API, dan tata kelola data pribadi?
\end{quote}

Kontribusi awal makalah dapat ditulis:

\begin{quote}
Makalah ini menawarkan analisis studi kasus Bank Wanua dengan memetakan
ketidakselarasan antara strategi bisnis dan kapabilitas arsitektur enterprise.
Kontribusi praktisnya adalah rancangan awal arah BDAT, prioritas API, dan
prinsip tata kelola yang realistis untuk bank daerah di Indonesia. Kontribusi
akademiknya adalah contoh penerapan TOGAF-lite untuk menjembatani strategi
UMKM digital, operating model, dan konteks regulasi lokal.
\end{quote}

Contoh ini belum sempurna, tetapi sudah cukup untuk menjadi dasar konsultasi.
Pada bab-bab berikutnya, RQ tersebut dapat dipersempit lagi sesuai data yang
tersedia dan target publikasi yang dipilih.

\section{Langkah Analisis dan Artefak Pendukung}

Gunakan delapan langkah berikut untuk mengubah topik awal menjadi rumusan
masalah, RQ, dan kontribusi awal.

\begin{enumerate}
  \item \textbf{Tulis ulang strategi organisasi dalam satu paragraf.} Hindari
        slogan. Sebutkan pelanggan, nilai, kanal, dan prioritas.
  \item \textbf{Petakan model bisnis secara ringkas.} Gunakan komponen kanvas
        model bisnis yang paling relevan. Tidak semua komponen harus sama
        panjang.
  \item \textbf{Tentukan operating model yang paling mendekati.} Apakah
        organisasi membutuhkan integrasi data tinggi, standardisasi proses
        tinggi, keduanya, atau justru otonomi unit?
  \item \textbf{Daftar ketidakselarasan utama.} Cari jarak antara strategi dan
        kondisi proses, data, aplikasi, teknologi, API, atau tata kelola.
  \item \textbf{Pilih satu ketidakselarasan utama.} Makalah satu semester
        sebaiknya tidak mencoba menjawab semua masalah sekaligus.
  \item \textbf{Tulis rumusan masalah 1--2 paragraf.} Jelaskan konteks,
        ketegangan, dampak, dan mengapa masalah ini penting dianalisis.
  \item \textbf{Ubah menjadi RQ.} Gunakan bentuk tanya yang dapat dijawab
        dengan analisis, bukan pertanyaan yang jawabannya sekadar ya atau
        tidak.
  \item \textbf{Tulis kontribusi awal.} Bedakan kontribusi praktis bagi
        organisasi dan kontribusi akademik bagi pembaca makalah.
\end{enumerate}

Tabel~\ref{tab:rq-baik-buruk} memberi contoh perbaikan RQ.

\begin{table}[htbp]
\centering
\scriptsize
\caption{Contoh memperbaiki pertanyaan penelitian.}
\label{tab:rq-baik-buruk}
\begin{tabularx}{\textwidth}{@{}p{0.30\textwidth} X X@{}}
\hline
\textbf{Versi Lemah} & \textbf{Masalahnya} & \textbf{Versi Lebih Kuat} \\
\hline
Bagaimana membuat aplikasi mobile Bank Wanua? & Terlalu teknis dan langsung
melompat ke solusi. & Bagaimana arsitektur aplikasi dan data Bank Wanua perlu
diselaraskan untuk mendukung layanan digital UMKM lintas kanal? \\
Apakah API penting bagi bank daerah? & Jawabannya terlalu umum dan mudah
ditebak. & Layanan API apa yang sebaiknya diprioritaskan bank daerah untuk
mendukung kemitraan UMKM, dan risiko tata kelola apa yang perlu dikendalikan? \\
Bagaimana transformasi digital rumah sakit? & Terlalu luas untuk satu makalah.
& Bagaimana rumah sakit daerah dapat menyelaraskan proses pendaftaran, data
pasien, dan kanal telekonsultasi dalam arsitektur target yang sesuai regulasi
data pribadi? \\
\hline
\end{tabularx}
\end{table}

\subsection*{Artefak Pendukung Bab Ini}

Bab ini menghasilkan artefak ringan, tetapi penting. Artefak tersebut tidak
dinilai sebagai luaran terpisah; fungsinya adalah memperkuat argumen awal
makalah.

\begin{table}[htbp]
\centering
\scriptsize
\caption{Artefak Bab 2 dan fungsi akademiknya.}
\label{tab:artefak-bab2}
\begin{tabularx}{\textwidth}{@{}p{0.26\textwidth} X X@{}}
\hline
\textbf{Artefak} & \textbf{Fungsi Analisis} & \textbf{Masuk ke Makalah Sebagai} \\
\hline
Ringkasan strategi & Menjelaskan arah organisasi dan prioritas transformasi. &
Latar belakang pada Pendahuluan. \\
Kanvas model bisnis ringkas & Menghubungkan pelanggan, nilai, aktivitas,
sumber daya, dan mitra. & Konteks organisasi atau lampiran ringkas. \\
Klasifikasi operating model & Menjelaskan kebutuhan standardisasi proses dan
integrasi data. & Justifikasi masalah dan metode analisis. \\
Daftar ketidakselarasan & Menunjukkan jarak antara strategi dan kapabilitas
saat ini. & Bukti awal rumusan masalah. \\
Draf RQ dan kontribusi & Mengarahkan struktur analisis semester berjalan. &
Bagian akhir Pendahuluan. \\
\hline
\end{tabularx}
\end{table}

Artefak yang baik harus relevan dengan RQ, mudah dipahami pembaca non-IT,
menunjukkan sebab-akibat atau prioritas, konsisten dengan narasi makalah, dan
menyebutkan sumber data atau asumsi yang digunakan.

\section{Integrasi ke Makalah Publikasi}

Bab ini terutama memperkuat bagian \textbf{Pendahuluan} makalah. Setelah
menyelesaikan bab ini, mahasiswa seharusnya memiliki empat bahan yang langsung
dapat dipindahkan ke naskah:

\begin{enumerate}
  \item satu paragraf konteks organisasi dan strategi digital;
  \item satu paragraf rumusan masalah;
  \item satu atau dua RQ yang dapat dijawab melalui analisis;
  \item satu paragraf kontribusi awal.
\end{enumerate}

Keempat bahan ini juga menjadi dasar bagian \textbf{Metode} dan
\textbf{Analisis} pada bab-bab berikutnya. RQ menentukan data apa yang perlu
dikumpulkan, kerangka apa yang dipakai, dan artefak apa yang harus dibuat. Jika
RQ menekankan integrasi data, Bab 5 akan menjadi sangat penting. Jika RQ
menekankan layanan API, Bab 9 dan Bab 10 perlu mendapat porsi lebih besar. Jika
RQ menekankan platform, Bab 11 dan Bab 12 akan menjadi inti diskusi.

Ingat bahwa kontribusi awal masih boleh berubah. Pada tahap ini, kontribusi
berfungsi sebagai arah kerja. Setelah literatur dibaca dan analisis kasus
dilakukan, kontribusi dapat dipertajam menjadi lebih spesifik.

\subsection*{Contoh Paragraf Akademik}

Contoh berikut menunjukkan cara mengubah artefak Bab 2 menjadi narasi
Pendahuluan:

\begin{quote}
Organisasi studi kasus menunjukkan ketidakselarasan antara ambisi digital dan
model operasi yang tersedia. Meskipun strategi organisasi menekankan layanan
digital lintas kanal, proses inti masih berjalan berbeda antarunit, sementara
data pelanggan belum memiliki definisi dan kepemilikan yang konsisten. Kondisi
ini menunjukkan bahwa masalah utama bukan sekadar keterbatasan aplikasi,
melainkan belum tersedianya arah arsitektur enterprise yang menghubungkan
strategi, integrasi data, standardisasi proses, dan tata kelola. Oleh karena
itu, penelitian ini mengajukan pertanyaan bagaimana pendekatan Enterprise
Architecture dapat digunakan untuk menyelaraskan strategi digital organisasi
dengan kapabilitas bisnis dan teknologi yang realistis dalam konteks Indonesia.
\end{quote}

\section{Diskusi Kritis dan Perangkap Umum}

Beberapa perangkap berikut sering terjadi ketika mahasiswa merumuskan masalah
penelitian.

\begin{enumerate}
  \item \textbf{Masalah terlalu luas.} ``Transformasi digital sektor kesehatan
        Indonesia'' tidak cocok untuk makalah satu semester. Persempit menjadi
        organisasi, proses, kapabilitas, atau keputusan tertentu.
  \item \textbf{Masalah terlalu teknis.} ``Implementasi REST API'' belum tentu
        menjadi masalah EA. Jelaskan nilai bisnis, mitra, tata kelola, dan
        risiko yang membuat API itu penting.
  \item \textbf{RQ tidak dapat dijawab dengan data yang tersedia.} Jika
        mahasiswa tidak punya akses ke data internal, pilih pendekatan yang
        mengandalkan dokumen publik, wawancara terbatas, atau kajian
        konseptual berbasis kasus.
  \item \textbf{Kontribusi terdengar seperti janji konsultan.} Kalimat seperti
        ``meningkatkan efisiensi 50\%'' perlu bukti. Untuk makalah awal,
        kontribusi yang lebih aman adalah kerangka analisis, rancangan awal,
        atau rekomendasi berbasis bukti.
  \item \textbf{Mengabaikan target publikasi.} Target publikasi memengaruhi
        panjang, gaya sitasi, jumlah referensi, dan kedalaman metode. RQ harus
        cukup cocok dengan ruang lingkup target publikasi.
  \item \textbf{Tidak membedakan gejala dan masalah.} Aplikasi lambat,
        data ganda, atau proses manual sering hanya gejala. Masalah yang lebih
        dalam mungkin terletak pada model operasi, kepemilikan data, atau
        prioritas strategi.
\end{enumerate}

Trade-off utama pada bab ini adalah antara keluasan dan ketajaman. Rumusan
yang terlalu sempit mudah dianalisis tetapi mungkin kurang bernilai akademik.
Rumusan yang terlalu luas terdengar penting tetapi sulit diselesaikan. Pilih
masalah yang cukup kecil untuk ditangani, tetapi cukup penting untuk
dipublikasikan.

\begin{table}[htbp]
\centering
\scriptsize
\caption{Trade-off saat merumuskan masalah dan RQ.}
\label{tab:tradeoff-rq}
\begin{tabularx}{\textwidth}{@{}p{0.22\textwidth} X X X@{}}
\hline
\textbf{Keputusan} & \textbf{Manfaat} & \textbf{Risiko} & \textbf{Prioritas} \\
\hline
RQ sangat sempit & Analisis lebih mudah dikendalikan. & Kontribusi dapat
terlihat terlalu kecil. & Cocok bila data terbatas dan target publikasi
praktis. \\
RQ luas lintas domain & Kontribusi terasa strategis. & Sulit selesai dalam
satu semester. & Hanya dipilih jika data dan literatur cukup kuat. \\
Fokus pada teknologi & Artefak teknis lebih jelas. & Argumen manajerial
melemah. & Perlu selalu dikaitkan dengan strategi dan risiko. \\
Fokus pada tata kelola & Relevan bagi manajemen dan regulasi. & Butuh bukti
tentang peran, proses, dan keputusan. & Cocok untuk organisasi yang datanya
sensitif atau bermitra luas. \\
\hline
\end{tabularx}
\end{table}

\section{Latihan Terapan Individual}

\subsection*{Latihan~2.1: Ringkasan Strategi dan Model Bisnis}

Gunakan organisasi studi kasus yang sudah dipilih pada Bab~\ref{ch:orientasi}.
Tuliskan ringkasan satu sampai dua halaman yang memuat:

\begin{enumerate}
  \item strategi digital organisasi dalam satu paragraf;
  \item tiga segmen pengguna atau pemangku kepentingan utama;
  \item proposisi nilai yang dijanjikan kepada masing-masing segmen;
  \item aktivitas kunci, sumber daya kunci, dan mitra kunci;
  \item satu masalah strategis yang paling tampak dari peta tersebut.
\end{enumerate}

\subsection*{Latihan~2.2: Klasifikasi Operating Model}

Nilailah organisasi Anda menggunakan dua pertanyaan:

\begin{enumerate}
  \item Seberapa tinggi kebutuhan standardisasi proses?
  \item Seberapa tinggi kebutuhan integrasi data?
\end{enumerate}

Pilih salah satu pola operating model: \emph{Diversification},
\emph{Coordination}, \emph{Replication}, atau \emph{Unification}. Jelaskan
alasan pemilihan tersebut dalam 2--3 paragraf. Jika organisasi Anda berada di
antara dua pola, jelaskan ketegangan yang muncul.

\subsection*{Latihan~2.3: Draf Rumusan Masalah, RQ, dan Kontribusi}

Susun draf awal dengan format berikut:

\begin{enumerate}
  \item Rumusan masalah: 1--2 paragraf.
  \item RQ utama: satu pertanyaan.
  \item Sub-RQ opsional: maksimum dua pertanyaan.
  \item Kontribusi praktis: 2--3 kalimat.
  \item Kontribusi akademik: 2--3 kalimat.
\end{enumerate}

Pastikan RQ tidak langsung melompat ke solusi. RQ yang baik membuka ruang
analisis, bukan sekadar meminta pembaca menyetujui teknologi tertentu.

\subsection*{Latihan~2.4: Uji Kelayakan Target Publikasi}

Ambil satu target publikasi dari daftar Bab~\ref{ch:orientasi}. Periksa apakah
rumusan masalah dan RQ Anda cocok dengan ruang lingkup target tersebut.
Tuliskan catatan singkat:

\begin{itemize}
  \item bagian ruang lingkup target publikasi yang paling relevan;
  \item tipe artikel yang paling cocok;
  \item panjang naskah atau batas halaman;
  \item konsekuensi bagi struktur RQ dan metode.
\end{itemize}

\subsection*{Rubrik Mini Latihan}

\begin{table}[htbp]
\centering
\scriptsize
\caption{Rubrik mini untuk keluaran latihan Bab 2.}
\label{tab:rubrik-latihan-bab2}
\begin{tabularx}{\textwidth}{@{}p{0.26\textwidth} X@{}}
\hline
\textbf{Kriteria} & \textbf{Indikator Kualitas} \\
\hline
Relevansi masalah & Masalah berasal dari strategi organisasi, bukan sekadar
ketertarikan pada teknologi tertentu. \\
Ketepatan konsep & Kanvas model bisnis dan operating model digunakan secara
tepat untuk menjelaskan ketidakselarasan. \\
Kejelasan artefak & Ringkasan strategi, klasifikasi operating model, dan
daftar ketidakselarasan mudah dibaca audiens non-IT. \\
Kekuatan argumen & Rumusan masalah menunjukkan konteks, bukti awal, dampak,
dan urgensi akademik/praktis. \\
Kelayakan implementasi & RQ dapat dijawab dengan data, waktu, dan kapasitas
yang tersedia dalam satu semester. \\
Keterhubungan ke makalah & Kontribusi awal jelas masuk ke bagian Pendahuluan
dan memberi arah bagi metode serta analisis. \\
\hline
\end{tabularx}
\end{table}

\subsection*{Pertanyaan Refleksi}

\begin{enumerate}
  \item Apakah masalah yang Anda pilih berasal dari strategi organisasi atau
        hanya dari ketertarikan pada teknologi tertentu?
  \item Apakah RQ Anda dapat dijawab dengan data dan waktu yang tersedia dalam
        satu semester?
  \item Apa risiko terbesar jika organisasi mengikuti rekomendasi yang mungkin
        muncul dari makalah Anda?
  \item Bagian mana yang perlu dikonsultasikan dengan dosen sebelum masuk ke
        Bab 3?
\end{enumerate}

\section{Ringkasan, Refleksi, dan Jembatan}

\subsection*{Poin Kunci}

\begin{itemize}
  \item Strategi bisnis harus dibaca sebagai pilihan tentang pelanggan,
        proposisi nilai, prioritas, dan cara organisasi menciptakan nilai.
  \item Kanvas model bisnis membantu menghubungkan strategi dengan aktivitas,
        sumber daya, mitra, kanal, dan struktur biaya.
  \item Operating model membantu membaca kebutuhan standardisasi proses dan
        integrasi data.
  \item Rumusan masalah yang baik menjelaskan ketegangan, bukti awal, dampak,
        dan konteks.
  \item RQ harus dapat dijawab melalui analisis, bukan sekadar menyatakan
        solusi yang sudah dipilih sejak awal.
  \item Kontribusi awal perlu membedakan nilai praktis bagi organisasi dan
        nilai akademik bagi pembaca makalah.
\end{itemize}

\subsection*{Jawaban Ringkas atas Pertanyaan Pemandu}

\begin{enumerate}
  \item Strategi organisasi dibaca dari pelanggan prioritas, nilai yang
        dijanjikan, kanal, aktivitas kunci, dan mitra kunci.
  \item Model bisnis menunjukkan bagaimana organisasi menghasilkan nilai,
        sedangkan operating model menunjukkan bagaimana nilai itu dijalankan
        melalui proses dan data.
  \item Ketidakselarasan muncul ketika strategi menuntut kapabilitas yang
        belum didukung proses, data, aplikasi, teknologi, API, atau tata
        kelola.
  \item Bukti awal dapat berupa dokumen publik, proses yang diamati, wawancara,
        laporan tahunan, regulasi, atau artefak internal yang diizinkan.
  \item Hasil analisis menjadi argumen akademik ketika artefak dipakai untuk
        mendukung klaim, bukan hanya ditempel sebagai gambar.
\end{enumerate}

\subsection*{Uji Mandiri: Pertanyaan Konseptual}

\begin{enumerate}
  \item Mengapa strategi bisnis perlu dibaca sebelum menyusun arsitektur
        enterprise?
  \item Jelaskan perbedaan antara strategi bisnis dan model bisnis.
  \item Apa fungsi operating model dalam analisis Enterprise Architecture?
  \item Sebutkan dua dimensi utama operating model dan jelaskan artinya.
  \item Apa perbedaan antara rumusan masalah, RQ, dan pernyataan kontribusi?
  \item Mengapa RQ yang terlalu teknis dapat melemahkan makalah Enterprise
        Architecture?
\end{enumerate}

\subsection*{Pertanyaan Aplikatif}

\begin{enumerate}
  \item Sebuah universitas ingin membangun layanan akademik terpadu. Mahasiswa,
        dosen, keuangan, dan administrasi memakai aplikasi berbeda. Rumuskan
        satu masalah penelitian yang lebih tajam daripada ``sistem belum
        terintegrasi''.
  \item Sebuah koperasi ingin menjual produk anggotanya melalui marketplace
        lokal. Pilih operating model yang paling mendekati dan jelaskan
        alasannya.
  \item Sebuah rumah sakit daerah ingin memperluas telekonsultasi, tetapi data
        pasien tersebar antara unit pendaftaran, poli, laboratorium, dan
        farmasi. Buat satu RQ yang menghubungkan strategi layanan digital dan
        arsitektur data.
\end{enumerate}

\subsection*{Pertanyaan Reflektif untuk Makalah}

\begin{enumerate}
  \item Apakah RQ Anda lebih dekat dengan domain Business, Data, Application,
        Technology, API, atau platform? Apa konsekuensinya bagi bab-bab
        berikutnya?
  \item Apakah kontribusi yang Anda tulis lebih kuat sebagai kontribusi
        praktis, akademik, atau keduanya? Mengapa?
\end{enumerate}

\subsection*{Ringkasan Naratif}

Bab ini menunjukkan bahwa rumusan masalah makalah tidak boleh muncul dari
ketertarikan teknologi semata. Rumusan yang baik berangkat dari strategi
organisasi, model bisnis, model operasi, dan ketidakselarasan yang dapat
dibuktikan. Kanvas model bisnis membantu pembaca melihat bagaimana organisasi
menciptakan nilai, sedangkan operating model membantu membaca kebutuhan
standardisasi proses dan integrasi data.

Dalam konteks Indonesia, strategi digital harus mempertimbangkan kematangan
organisasi yang beragam, regulasi seperti UU PDP dan PP No.~71/2019,
ketentuan sektor seperti SNAP BI, struktur organisasi yang sering hierarkis,
serta keterbatasan biaya dan SDM. Kepekaan terhadap konteks ini membuat RQ
lebih realistis dan kontribusi lebih dapat dipercaya.

Keluaran utama bab ini adalah draf rumusan masalah, RQ, dan kontribusi awal.
Ketiganya akan menjadi kompas bagi bab-bab berikutnya. Jika kompas ini kabur,
artefak arsitektur yang dibuat pada Bab 3 sampai Bab 12 akan mudah menjadi
kumpulan gambar yang tidak menyatu dengan argumen makalah.

\subsection*{Jembatan ke Bab Berikutnya}

Bab berikutnya memperkenalkan TOGAF-lite dan cara menyusun tinjauan pustaka
awal. Setelah strategi dan RQ mulai jelas, pembaca perlu memilih kerangka
analisis yang dapat dipertanggungjawabkan. TOGAF-lite akan dipakai sebagai
bahasa kerja untuk menyusun visi arsitektur, kondisi saat ini, arsitektur
target, analisis kesenjangan, dan peta jalan. Pada saat yang sama, mahasiswa
mulai membangun bibliografi awal agar rumusan masalah tidak hanya kuat secara
praktis, tetapi juga tersambung dengan percakapan akademik.

\section{Bacaan Lanjutan}

\subsection*{Bacaan Inti}

\begin{itemize}
  \item \emph{Enterprise Architecture as Strategy}~\cite{ross2006enterprise}.
        Rujukan utama untuk memahami operating model dan hubungan strategi
        bisnis dengan fondasi arsitektur.
  \item \emph{Designed for Digital}~\cite{ross2019designed}. Membantu pembaca
        memahami kapabilitas organisasi digital dan mengapa strategi perlu
        diterjemahkan menjadi kemampuan yang dapat dijalankan.
  \item \emph{TOGAF Standard, 10th Edition}~\cite{togaf10}. Bacaan awal untuk
        melihat bagaimana visi arsitektur dan kebutuhan bisnis ditempatkan
        dalam siklus pengembangan arsitektur.
\end{itemize}

\subsection*{Bacaan untuk Memperdalam}

\begin{itemize}
  \item \emph{Platform Revolution}~\cite{parker2016platform}. Berguna bagi
        mahasiswa yang rumusan masalahnya mengarah ke ekosistem platform dan
        \emph{network effects}.
  \item \emph{The Business of Platforms}~\cite{cusumano2019business}. Cocok
        untuk memperdalam diskusi tentang strategi dan tata kelola model
        bisnis platform.
  \item \emph{APIs: A Strategy Guide}~\cite{jacobson2011apis}. Bermanfaat jika
        ketidakselarasan utama organisasi berkaitan dengan kemitraan API,
        integrasi layanan, atau pembukaan kanal digital untuk pihak luar.
\end{itemize}

\subsection*{Sumber Indonesia}

\begin{itemize}
  \item \emph{Undang-Undang Pelindungan Data Pribadi}~\cite{uupdp2022}.
        Penting untuk mahasiswa yang rumusan masalahnya menyentuh data
        pelanggan, personalisasi layanan, atau integrasi data lintas kanal.
  \item \emph{PP No.~71/2019}~\cite{pp712019}. Relevan untuk memahami
        tanggung jawab penyelenggara sistem elektronik.
  \item \emph{Standar Nasional Open API Pembayaran (SNAP BI)}~\cite{snapbi}.
        Sangat relevan untuk topik open banking, fintech, pembayaran digital,
        dan API sektor keuangan.
\end{itemize}

\subsection*{Lampiran: Worksheet Ringkas Bab 2}

Gunakan format berikut sebagai catatan kerja sebelum menulis bagian
Pendahuluan makalah.

\begin{enumerate}
  \item \textbf{Strategi organisasi:} tuliskan 4--6 kalimat tentang pelanggan,
        nilai, kanal, prioritas, dan batasan utama.
  \item \textbf{Model bisnis:} rangkum segmen pelanggan, proposisi nilai,
        aktivitas kunci, sumber daya kunci, dan mitra kunci.
  \item \textbf{Operating model:} pilih salah satu pola dan jelaskan alasan
        berdasarkan kebutuhan standardisasi proses dan integrasi data.
  \item \textbf{Ketidakselarasan utama:} tuliskan tiga jarak paling penting
        antara strategi dan kondisi saat ini.
  \item \textbf{Rumusan masalah:} tulis 1--2 paragraf yang menjelaskan konteks,
        ketegangan, dampak, dan pentingnya analisis.
  \item \textbf{RQ dan kontribusi:} tulis satu RQ utama, maksimum dua sub-RQ,
        kontribusi praktis, dan kontribusi akademik.
\end{enumerate}
